\quad O Bubble Sort é um algoritmo de ordenação baseado na comparação
e troca de elementos adjacentes. Seu funcionamento consiste em percorrer
o vetor várias vezes, comparando elementos vizinhos e realizando uma troca
quando eles estão fora da ordem desejada. Dessa forma, os maiores elementos
vão sendo deslocados gradualmente para o final do vetor, de maneira
semelhante ao movimento de bolhas subindo em um líquido.

\quad A cada passagem pelo vetor, os elementos adjacentes são comparados.
Caso o elemento da esquerda seja maior que o elemento da direita, suas
posições são trocadas. Após cada passagem, pelo menos um dos maiores
elementos ainda não ordenados é colocado em sua posição definitiva no final
do vetor. O processo continua até que todos os elementos estejam
ordenados.

\begin{lstlisting}[caption={Implementação do Bubble Sort}]
void bubbleSort(int* vetor, int tamanho) {
    for (int i = 0; i < tamanho - 1; i++) {
        bool trocou = false;

        for (int j = 0; j < tamanho - i - 1; j++) {
            if (vetor[j] > vetor[j + 1]) {
                int temp = vetor[j];
                vetor[j] = vetor[j + 1];
                vetor[j + 1] = temp;
                trocou = true;
            }
        }
        if (!trocou)
            break;
    }
}
\end{lstlisting}

\quad No primeiro laço, a variável $i$ representa a quantidade de
passagens já realizadas pelo vetor. O segundo laço percorre os elementos
ainda não ordenados, comparando cada elemento com seu vizinho à direita.

\quad Quando o elemento da esquerda é maior que o elemento da direita,
é realizada uma troca entre eles utilizando a variável auxiliar
\texttt{temp}. Ao final de cada passagem, o maior elemento da parte ainda
não ordenada é levado para o final do vetor. Por isso, o limite do segundo
laço diminui a cada nova passagem.
